% Folder 74, batch 5 (archivists' pages 81-100).
% Pass: Opus 5.5 (claude-opus-5-5), 2026-10-03 — first pass, unchecked against the pages by a human.
%
% Fast drafting hand throughout: formulas carry, connective prose often does not.
% Notation fixed for the batch: X a "pavé" (the graph of a 3x3 rook's board,
% i.e. the product of two triangles), C a square ("carré") in X, a its centre;
% Car_?(X), Hex_?(X) for the sets of squares / hexagons with an extra structure
% (subscripts d, s, *, ép, tr as on the page); \Gamma^\circ the opposite
% (complementary) graph; Pav(C) the pavé enveloping C.
\documentclass[11pt,a4paper]{article}
\input{../preamble/grothendieck}

\folder{74}
\batch{5}
\pages{81}{100}
\foldertitle{Complexes cubiques : notes manuscrites (s.d.).}
\dating{[à partir de 1976]}
\watermark{Édition de démonstration}

\begin{document}

\section{Carrés et hexagones dans un pavé}

\page{81}
\emph{Carrés et hexagones dans un pavé.}

$X$ pavé.

I) L'ens des carrés $C\subset X$ est en corr $1$-$1$ avec l'ens (des sommets de) $X$, à tout carré $C$ correspondant le \uncertain{sommet} unique \add{et : $a\notin C$, l'un des pts de $C$ non liés : \ill{}} $x\in X$ non lié \struck{aux sommets} \add{à tous les sommets} de $C$, [$C$ \emph{NB} il y a un seul pt de $X$ \add{non} lié à \struck{deux pt} deux pts distincts non liés --- i.e.\ à \struck{\uncertain{angle}} \add{paire} de pts \struck{distincts} non liés est contenue dans une triade unique --- la donnée d'un tel couple équivaut \add{dans $C$} à la donnée d'une triade pointée $(t,a)$, en associant à $(t,a)$ la paire $t\setminus\{a\}$~]. Pour que 2 paires de pts non liés $\{x,x'\}$, $\{y,y'\}$ forment un carré, il faut et il suffit que \struck{les triades} leurs ``centres'' soient les mêmes, i.e.\ qu'elles soient disjointes et \struck{\uncertain{orthogonales}} : leurs triades \emph{sécantes} ; la donnée d'une \struck{diagonale} des paires $\{x,x'\}$ détermine l'autre sans ambiguïté ; $\{y,y'\}$ est aussi l'ens des pts de $X$ liés à aucun des $\{x,x'\}$. Donc la donnée d'\uncertain{un carré à diagonale marquée} équivaut à celle d'une paire de pts non liés, i.e.\ encore à celle d'une triade pointée.

La donnée d'un \uncertain{fanion} \ill{} équivaut à celle d'une triade \uncertain{épinglée}.
\[
  \mathrm{Car}_{d}(X)\simeq \mathrm{tria\,pt}(X),\qquad
  \mathrm{Car}_{s}(X)\simeq \mathrm{tria\,ép}(X)
\]
(\uncertain{est un $\mathfrak{S}_3$-\ill{} cubique}).
\note{Devant $\mathrm{Car}_d$, un mot biffé illisible ; les indices $d$, $s$ sont lus d'après la page suivante, où les mêmes symboles reviennent.}

\add{Chacun des 4} \struck{Les autres} pts de $X\setminus C$ distincts du centre $a$ de $C$ est lié exactement à \uncertain{un point} de $C$ diag $\{x,x'\}$, du …
\drawing{Dans la marge gauche, des schémas de points (gros points, petits points, croix) et de segments pleins ou pointillés : un segment et une demi-diagonale pointillée en haut, puis deux cadres de $3\times 3$ points reliés par des flèches, illustrant la position d'un carré et de son centre dans le pavé.}

\page{82}
du \ill{} par l'autre diagonale $\{y,y'\}$, \ill{} $Y=X\setminus(C\cup\{a\})$ s'envoie dans le produit des diagonales et on trouve là une bijection. Les 4 pts sont d'ailleurs liés à $a$, \uncertain{ainsi que} les 4 sommets des deux triangles de sommet $a$. D'où \uncertain{aussi} que l'\uncertain{ombre} sur $C$ d'un pt de $Y$ est un côté de $C$, et qu'on trouve une corr.\ $1$-$1$ de $C$ sur l'un des 4 côtés de $C$. \struck{Donc} Deux pts de $Y$ sont liés ssi \struck{$\mathrm{Car}_{s*}(C)\simeq$} \struck{Alors} ils correspondent à des côtés opposés de $C$. Donc les 2 paires de côtés opposés de $C$ correspondent aux 2 triangles de sommet $a$.

Donc
\begin{align*}
  \mathrm{Car}_{d*}(X) &\simeq \mathrm{tria\,pt}(X),\\
  \mathrm{Car}_{s*}(X) &\simeq \mathrm{arc}(X)\simeq \mathrm{tria\,ép}(X)
\end{align*}
(\uncertain{ici un $\mathfrak{S}_3$-torseur libre}~!)
\note{À gauche de ces formules, des schémas : un triangle à un sommet marqué ; une flèche $\simeq$ un triangle à arête fléchée ; dans de petits ovales, des légendes de traits (plein, gras, pointillé).}

D'ici
\[
  \mathrm{Car}_{\cdot}(X)\simeq \text{ens des couples d'une triade et deux triangles sécants dans } X.
\]
\drawing{Un triangle plein de sommet marqué, flanqué de deux segments pointillés.}

Enfin
$\mathrm{Car}_{r}(X)\simeq{}$ \uncertain{épinglages des} pavé
\drawing{Un triangle de sommet marqué, une croix et un point, « ou » une flèche verticale avec un segment pointillé.}

alors
$\mathrm{Car}_{\omega}(X) ={}$ \struck{\ill{}} $\mathrm{Car}(X)/\mathrm{opp}$
(\uncertain{couples} d'une triade et deux triangles sécants, aux couples opposés).
\drawing{Dans la marge, un schéma biffé de hachures.}

\drawing{Grand diagramme écrit en travers de la page (feuille tournée), reliant des ensembles d'hexagones : $\mathrm{Hex}_{s}\simeq\mathrm{hex}$ ($\simeq$ tria), $\mathrm{Hex}_{\mathrm{tr}}$ ($\simeq$ \ill{} tria), $\mathrm{Hex}_{s*}$, $\mathrm{Hex}_{\mathrm{tr}*}$ (triplets de \ill{}), $\mathrm{Hex}_{\mathrm{ép}}$, $\mathrm{Hex}_{\mathrm{tri}}$, $\mathrm{Hex}_{d}$ ($\simeq$ tria \ill{}), $\mathrm{Hex}$ ($\simeq$ pav) et « tri hex », par des flèches horizontales marquées $3$ et des flèches verticales et obliques marquées $2$, formant un réseau en forme de prisme.}
\note{Le diagramme est trop serré pour être reporté en \texttt{tikz-cd} sans deviner le sens de plusieurs flèches ; il est décrit.}

\page{83}
exactement --- \ill{} correspondant : 4 carrés de \ill{} centre $a$~].

\emph{Question} Quels sont les graphes finis qui ont une ``enveloppe'' \add{pavée} \uncertain{trilatère}~? Exemples.
\drawing{Rangée de graphes : un point, un segment, un carré, un hexagone, un pentagone, deux segments (dans un ovale), chacun avec une flèche vers le bas ; sous le dernier, un sablier (bitriangle). À gauche, entouré, un hexagone contenant un carré et une croix.}
(et leurs sommes, \add{pavés} \uncertain{trilatères} mais $\neq\emptyset$…)

\struck{\ill{}} \uncertain{dessus} pavé carré … et les \emph{\uncertain{quadrilatères}} \ill{} et leurs enveloppes.

Regardons p.\ ex.
\drawing{Cinq esquisses de pyramides et de bipyramides en perspective (sommets marqués, arêtes croisées), dont deux barrées d'une croix au crayon, et une entourée d'un arc ; à côté, des pointillés.}
\struck{\ill{} Regardons}

\drawing{Sous un trait horizontal : un petit carré « pavé », un cercle « \uncertain{coque} », un double ovale « \uncertain{haltère} ».}

\drawing{Diagramme esquissé : $\mathrm{Car}_{s}=\mathrm{hex}_{d}\longrightarrow\mathrm{Car}_{s,\mathrm{ép}}$, avec $\wr$ vers $\mathrm{tria\,ép}$ et $\mathrm{Hex\,ép}$ ; puis $\mathrm{Car}_{\wedge}\simeq\mathrm{hex}_{d}\longleftarrow\mathrm{Car}_{p'}\simeq\mathrm{hex}_{s*}$, $\wr\ \simeq\mathrm{hex}^{*}_{d}$, « \uncertain{tri pav.\ pet} » ; $\mathrm{Car}\simeq$ \uncertain{pav.\ pet.} ; $\mathrm{Car}_{d*}$, $\mathrm{Car}_{s*}$ ; chaque terme accompagné de petits schémas (triangles, segments pointillés, cadres).}

\drawing{En bas de page, un petit carré de flèches : $\mathrm{Car}_{d}\to\mathrm{Car}$ (lettres soulignées deux fois) au-dessus de $\mathrm{hex}_{d}\to\mathrm{hex}$, les verticales marquées $\wr$, et une flèche oblique partant de $\mathrm{Car}$.}

\page{84}
\emph{NB} \struck{La correspondance} $C\mapsto X\setminus\{C\cup\{\cdot\}\}$ \struck{est}

\add{Appelons} \emph{anti-carré} l'\uncertain{image} d'un \struck{carrés de $X$ et} \uncertain{\ill{}}, i.e.\ un graphe somme de \uncertain{quatre} segments. Les anti-carrés dans $X$ sont des $X^{\circ}$, puis on a \uncertain{pour} \ill{} associés aux pts de $X^{\circ}$, \ill{} $X^{\circ}$ l'un des 4 pts de $X^{\circ}$ \uncertain{isomorphes} : $a$, i.e.\ que \ill{} $X^{\circ}$ liés à $a$. \struck{\ill{} les D…}

La correspondance $C\mapsto \bigl(X\setminus C\cup\{\text{centre de }C\}\bigr)$ entre carrés de $X$ et anti-carrés de $X$ est bijective. Si $C$ et $C'$ sont associés ainsi, $C$ et $C'$ sont des carrés \uncertain{qui ont} \uncertain{les sommets} de l'un étant aux côtés de l'autre.

Il y a aussi \uncertain{une corr.} \uncertain{ponctuelle} \add{(via $\mathrm{Pav}(C)$)} associant à chaque carré $C$ une \uncertain{paire} (la \uncertain{centrale} ``enveloppe'') dans l'un des sommets de $C$~: $C\amalg C^{*\circ}\amalg\{a\}$. Des \uncertain{points} de $C$ ($C^{*\circ}$) sont liés seulement à \ill{} de $C$ (\add{les} liés des \ill{} $C^{*\circ}$, i.e.\ les liés dans $C^{*}$), et un pt de $C$ est lié à un pt $x^{*}$ de $C^{*\circ}$ \struck{\ill{}} \ill{} $x$, $x^{*}$ incidents, \ill{} une \uncertain{pté} de $C$ liée aux

\page{85}
… de $C^{*\circ}$. L'\uncertain{union} de ce pavé est can.\ isomorphe \struck{\ill{}} à $\mathrm{Pav}(C^{*})$ \ill{} \struck{\ill{}} \add{utilisant} \struck{\ill{}} $C^{**}\simeq C$ --- qui est \ill{} avec l'antipodisme …
\[
  \mathrm{Pav}(C)^{\circ} \;\overset{\mathrm{canon}}{\simeq}\; \mathrm{Pav}(C^{*})
\]
\note{Sous le membre de gauche, « \uncertain{carré donné} » ; sous le membre de droite, « \uncertain{carré dual} de $C$ ».}

\emph{NB} La \uncertain{paire} des \ill{} de $X$ \uncertain{en} relation avec triangles issus de $a$ et \add{\uncertain{fixant}} celui qui \struck{\uncertain{échange}} \add{\ill{}} l'antipodisme, une \ill{} \struck{\ill{}} $C$ et son $C^{*\circ}$.

$\mathrm{Car\,Pav}(C)$ a la propriété universelle évidente pour les complexes trilatères $\mathcal{X}$ :
\[
  \mathrm{Pl}(\mathrm{Pav}(C),\mathcal{X}) \xrightarrow{\ \sim\ } \mathrm{Pl}(C,\mathcal{X})
  \qquad\text{(plongements)}
\]
Si $\mathcal{X}$ \uncertain{est lui-même un pavé} :
\[
  \mathrm{Iso}(\mathrm{Pav}(C),\mathcal{X}) \simeq \mathrm{Pl}(C,\mathcal{X}).
\]

Donc $\mathcal{X}$ complexe trilatère donné ; \struck{\ill{}} la donnée d'un carré de $\mathcal{X}$ revient à la donnée d'un \emph{pavé pointé} \ill{} (\uncertain{ce qui revient} \ill{}…)

\emph{NB} On \uncertain{peut} \ill{} que dans un graphe \emph{anticarré} (\uncertain{simplement} \ill{} des bitriangles $\bowtie$) est contenu dans \emph{4} pavés…

\page{86}
graphe produit de graphes $(\Gamma_i)_{i\in I}$ : ens des sommets $\prod\Gamma_i$ \struck{\uncertain{distincts}}, $x=(x_i)$ et $y=(y_i)$ liés ssi $\exists\, i\in I$ tel que $x_j=y_j$ si $j\neq i$, et $x_i,y_i$ liés (\ill{} \add{(par produit catégorique \ill{})}).

graphe \struck{\ill{}} \emph{opposé} $\Gamma^{\circ}$ du graphe $\Gamma$ : \uncertain{mêmes} sommets, deux pts \add{distincts} de $\Gamma^{\circ}$ sont liés ssi ils ne sont pas liés dans $\Gamma$. \uncertain{NB} pour une \uncertain{structure} \ill{} donnée, l'ens des \uncertain{structures} \ill{} \struck{$\Gamma$ correspondance bijective \ill{} \ill{} l'un des} $\mathrm{Ar}(\Gamma)$, le \uncertain{passage} \ill{} complémentaire. C'est une \uncertain{involution} \ill{}, \uncertain{pour} $\mathrm{Ar}(\Gamma)\neq\emptyset$ i.e.\ $\mathrm{card}(\Gamma)\geq 2$, sans pt fixe. Néanmoins, il \uncertain{vaut} \uncertain{mieux} que $\Gamma$ soit isomorphe au graphe \struck{\ill{}}. On dit alors que $\Gamma$ est ``\emph{auto-opposé}''.

\emph{Coproduit} de graphes $\Gamma_i$ : \uncertain{défini par} \add{graphe} $\Gamma=\overset{\circ}{\coprod}\Gamma_i$ (\uncertain{noté} $\Gamma_1\overset{\circ}{\times}\Gamma_2$ pour deux facteurs)
\[
  \Gamma^{\circ}=\prod\Gamma_i^{\circ},\quad\text{i.e.}\quad \Gamma=\Bigl(\prod\Gamma_i^{\circ}\Bigr)^{\circ}.
\]
Donc ens des sommets de $\Gamma=\prod\Gamma_i$. Deux \add{de $\Gamma$} \uncertain{distincts} $x=(x_i)$ et $y=(y_i)$ sont \emph{non} \add{liés dans $\Gamma$} (liés dans $\Gamma^{\circ}$), ssi $\exists\, i\in I$ tel que $x_j=y_j$ pour $j\neq i$, et $x_i,y_i$ \emph{non} liés dans $\Gamma_i$, i.e.\ $x_i,y_i$ liés dans $\Gamma_i^{\circ}$.
\marginal{\emph{NB} $\bigl(\prod\Gamma_i\bigr)^{\circ}=\overset{\circ}{\prod}\Gamma_i^{\circ}$ \uncertain{donc}.}

\emph{Prop} Un pavé (i.e.\ \add{graphe (\uncertain{sommets})} \uncertain{produit} \uncertain{de deux triangles}) est autodual, \uncertain{en d'autres termes} un pavé est isomorphe au coproduit de deux triades (duals des triangles) ou encore un \ill{} coproduit de deux triades est un pavé~?

\page{87}
Coproduit \add{d'un graphe} \struck{\ill{}} un doublet $\bullet\ \bullet$ donne, pour un graphe discret, les graphes bi-discrets.

Coproduit de deux doublets $(\overset{\bullet}{a}\ \overset{\bullet}{b})$, $(a'\ b')$ :
\drawing{Les quatre sommets $(ba')$, $(bb')$, $(aa')$, $(ab')$ disposés en carré, reliés par les deux diagonales qui se croisent.}

Coproduit des deux segments $(a\text{---}b)$, $(a'\text{---}b')$ :

\begin{tikzcd}
  (ba') \arrow[r, no head] \arrow[d, no head] & (bb') \arrow[d, no head] \\
  (aa') \arrow[r, no head] & (ab')
\end{tikzcd}

(1) carré.
\note{Entre les deux colonnes, un croisement de diagonales lourdement biffé.}

\uncertain{Idem} des précédents. [Le produit d'un doublet et d'un segment \struck{\ill{}} \add{\uncertain{disjoints}} \struck{segment $a'\text{---}b'$} est l'un d'\emph{entre} eux. Le coproduit de deux segments …{} le \uncertain{graphe} \uncertain{complété} \uncertain{semble} \ill{} …] Produit d'un doublet $(a\bullet\ \bullet b)$ et d'un segment $a'\text{---}b'$ :
\[
  (ba')\text{---}(bb'),\qquad (aa')\text{---}(ab'),
\]
son \uncertain{opposé} est le coproduit du segment $(a\text{---}b)$ et du doublet $(a'\cdot b')$ :
\drawing{(2) Un bitriangle (sablier couché) sur les sommets $(ba')$, $(bb')$, $(aa')$, $(ab')$ ; à côté, souligné : « carré ».}

Donc dans le produit de deux \uncertain{ensembles} de cardinal 2, on trouve 3 structures de carré naturelles distinctes, \struck{\ill{}} \uncertain{2 \ill{}}, \uncertain{produit} de deux segments…

\page{88}
\add{deux} et les structures \emph{co}produit obtenues en mettant sur l'un des facteurs une structure de doublet, sur l'autre une structure de segment. (Celle \uncertain{qu'on a} \emph{\uncertain{induite}} \uncertain{dessinée} est
\drawing{Un bitriangle sur les sommets $(ba')$, $(bb')$, $(aa')$, $(ab')$ : côtés horizontaux $(ba')\text{---}(bb')$ et $(aa')\text{---}(ab')$, une diagonale $(bb')\text{---}(aa')$ en trait fin, deux amorces de l'autre diagonale.}
On trouve ainsi \emph{les} trois \add{uniques} structures de carré sur l'ens à 4 él.\ produit, correspondant aux partitions de cet ens du type $(2,2)$…

Soit un carré $C$, défini par l'un (ou \uncertain{ses}) sommets, et la partition du type $2,2$ \add{(i.e.\ une involution sans pt fixe)} en les deux diagonales, $(C_{\alpha})_{\alpha\in D}$. \uncertain{Orientation} du carré définie, pour chaque $\alpha\in D$, \uncertain{par} une bijection $\varphi_{\alpha}: z\mapsto uz$ de $C_{\alpha}=\{x,x'\}$ sur $C_{\beta}=\{y,y'\}$.
\drawing{Un carré de sommets $x$, $y$ (en haut), $y'$, $x'$ (en bas), avec ses deux diagonales ; une flèche $u$ va de $x$ à $y$.}
On a
\[
  \varphi_{\beta}\varphi_{\alpha}=-\mathrm{id}_{C_{\alpha}},\qquad \varphi_{\alpha}\varphi_{\beta}=-\mathrm{id}_{C_{\beta}}
\]
(l'une de ces relations implique l'autre). D'où
\[
  \mathrm{or}(C)\simeq\Bigl(\underbrace{\bigwedge_{\alpha\in D}C_{\alpha}}_{\substack{D^{*}=\text{ens des diagonales}\\ \text{du carré dual}}}\Bigr)\times D \;\simeq\; D^{*}\times D .
\]
\note{Le symbole devant $C_\alpha$ est lu $\bigwedge$ ; « carré dual » sous l'accolade est une lecture incertaine.}

\page{89}
\note{Page au crayon, en largeur, occupée presque entièrement par un grand diagramme d'ensembles d'hexagones et de carrés du pavé, reliés par des flèches marquées $2$ ou $3$ (le degré des applications), avec en annotation le groupe opérant ($\mathfrak{S}_3$, $\mathfrak{S}_2$, $\mathbb{D}_3$, $\mathbb{D}_6$, $\mathbb{F}_2\times\mathbb{F}_2$). Elle reprend et étend les diagrammes des pp.~82 et~83.}

\marginal{invariants cube et hypercube dans domaines \struck{\ill{}} polyèdres et polytopes pro…{} (prismes $\left\{\begin{array}{l}\text{triangulaires}\\ \text{pentagonaux}\\ \text{hexagonaux}\end{array}\right.$ produit de polygones…)}

\drawing{Le diagramme. On y lit, sans pouvoir assurer le sens de toutes les flèches : en haut, $\mathrm{Car}_{d*}\simeq\mathrm{pav\,car}$, relié par des flèches marquées $2$ à $\mathrm{car}$, $\mathrm{car}_{\Omega}$ et $\mathrm{car}_{s*}$ ; au centre $S$, où aboutissent des flèches venant de $\mathrm{Ar}\xleftarrow{3}\mathrm{Arc}$ ($\simeq_3$ \uncertain{trig pt}, $\simeq_2$ trig ép), de « dodéc ép $\simeq$ pent ép $\simeq$ cour ép » et d'un schéma de « pent. » ; puis $\mathrm{hex}_{\mathrm{tr}*}$ ($\simeq$ pav pt $+$ tria), $\mathrm{hex}_{s*}\simeq\mathrm{car}_{r}$ ($\simeq$ pav $+$ tria $+$ \ill{}), $\mathrm{hex}_{\Omega^*}$ ($\mathfrak{S}_3$), $\mathrm{hex}_{\Omega}$ ($\mathfrak{S}_3$), $\mathrm{hex\,ép}\simeq\mathrm{car\,ép}$ ($G_{\mathrm{pav}}$ ; $\simeq$ pav ép, $\simeq$ car ép), $\mathrm{hex\,ép}'$ ($\mathfrak{S}_3$), $\mathrm{hex}_{d}$ ($\mathfrak{S}_2$ ; $\simeq$ tria pt $\simeq\mathrm{car}_{d}$) ; en bas, $\mathrm{Hex}$ ($\simeq$ tria \ill{}) $\xleftarrow[\mathrm{orr}]{2}\mathrm{Hex}_{s}\simeq\mathrm{hex}$ ($\simeq$ tria), $\mathrm{Hex}$ ($\simeq$ pav) $\xleftarrow[\mathrm{orr}]{2}$ ($\mathbb{D}_3$) $\mathrm{Hex}_{\mathrm{tr}}$ ($\simeq$ tr tria), ($\mathbb{D}_3$) $\mathrm{Hex\,ép}'$, ($\mathbb{D}_6$) $\mathrm{Hex\,ép}\simeq\mathrm{hex}_{\mathrm{tr}}$ ($\simeq$ tr ép tria ; $\mathfrak{S}_3$) $\xleftarrow{3}\mathrm{hex}_{s}\simeq\mathrm{car}_{s}$ ($\simeq$ tria ép ; $\mathfrak{S}_3$), $\mathrm{Hex}_{s*}$ ($\simeq$ ba pt), $\mathrm{Hex}_{\Omega}$, $\mathrm{Hex}_{t*}$ ($\simeq$ tri b tr), $\mathrm{Hex}_{tt*}$ ($\simeq$ tr $\circ$ tria ; $\mathbb{F}_2\times\mathbb{F}_2$), et la dernière ligne $\mathrm{triHex}$ ($\simeq$ tri pav) $\xleftarrow{3}\mathrm{pav}$, avec $\mathrm{Hex}\xrightarrow{3}\mathrm{triHex}$. Une bulle « échanges~! » renvoie par deux arcs de $\mathrm{hex\,ép}$ à lui-même. À gauche, une flèche « $\mathrm{b}\simeq$ doublet » et « bi pt ». À droite, des esquisses : trois prismes accolés, un pentagone subdivisé, une chaîne de points fléchés.}

\page{90}
\note{Au crayon, en largeur : mise au net partielle du diagramme de la p.~89.}

\drawing{Diagramme en forme de cube, flèches marquées $2$ ou $3$. En haut : $\mathrm{Car}_{d*}\simeq$ pav à \uncertain{trig pt}, recevant des flèches $2$ de $\mathrm{car}$ et de $\mathrm{Car}_{s*}\simeq$ pav.\ \uncertain{car} $\simeq$ pav $+$ trig ép ($\mathfrak{S}_3$). Au milieu : $\mathrm{hex}_{\Omega}\simeq$ tria $\Omega$ $\xleftarrow{3}\mathrm{hex}_{r'}\simeq\mathrm{car}_{r'}\simeq$ tria ép ($\mathfrak{S}_3$) ; $\mathrm{Car}\simeq$ pav pt. Puis $\mathrm{hex}$ ($\simeq$ tria) $\leftarrow\mathrm{hex}_{\mathrm{tr}*}$ ($\mathfrak{S}_3$) $\leftarrow\mathrm{hex}_{\mathrm{tr}\Omega}$ ($\simeq$ trig ép tria $\Omega$) $\xleftarrow{3}\mathrm{hex\,ép}$ ($\simeq$ trig ép tria ép, $\simeq$ pav ép ; $\mathfrak{S}_2$) ; $\mathrm{hex}_{d}$ ($\simeq\mathrm{Car}_{d}\simeq$ tria pt) $\leftarrow\mathrm{hex}_{s*}\leftarrow\mathrm{hex\,ép}$. En bas : $\mathrm{Pav\,coor}\simeq$ trig tria $\xleftarrow{3}\mathrm{hex}_{\mathrm{tr}}$ ($\mathfrak{S}_3$ ; $\simeq$ trig ép tria) $\xleftarrow{3}\mathrm{hex}_{s}$ ($\simeq$ tria ép ; $\mathfrak{S}_3$), trig tria pt $\simeq$ tria ; $\mathrm{Pav\,coor}\xrightarrow{2}\mathrm{Pav}$ ; trig $\Omega$ tria $\xleftarrow{3}$ tria $\Omega$. Les verticales $\mathrm{hex}\xrightarrow{3}\mathrm{Pav\,coor}$, $\mathrm{hex}_{\mathrm{tr}\Omega}\xrightarrow{2}\mathrm{hex}_{\mathrm{tr}}$, $\mathrm{hex}_{d}\xrightarrow{3}$ trig tria pt, $\mathrm{hex}_{s}\to$ tria $\Omega$ ; diagonales marquées $2$.}

\section{Complexes cubiques généralisés}
\note{La p.~91 est une chemise de couleur portant seulement, au crayon, « (33~p) » ; la liasse qu'elle annonce commence à la p.~92 et se poursuit au-delà de ce lot.}

\page{92}
\emph{Complexes cubiques généralisés} (complexes \emph{triadiques} ou \emph{trilatères}~?)

$S$ graphe strict, $A\subset\mathfrak{P}_2(S)$ l'ens des arêtes. On \struck{\ill{}} appelle \emph{triangle} de $S$ tout $t\in\mathfrak{P}_3(A)$ \note{lire sans doute $\mathfrak{P}_3(S)$} tel que $x,y\in t$, $x\neq y$ implique que $x,y$ sont ``liés'' (i.e.\ $\{x,y\}\in A$). Soit $T$ l'ens des triangles de $S$.

\emph{Axiome 1} $\forall\, a\in A$, $\exists!\, t\in T$ tel que $a\subset t$.

\emph{Axiome 2} $\forall\, t\in T$, \struck{\ill{}} $s\in S\setminus t$, $s$ est lié à un élément de $t$ (unique d'après l'axiome~1).

\emph{NB} La conjonction de 1 et 2 s'énonce aussi :

\emph{Axiome (C)} $\left\{\begin{array}{l}\text{a) toute arête est contenue dans un triangle}\\ \text{b) } \forall\, t\in T,\ s\in S\setminus t,\ \exists!\, s'\in t \text{ tel que } s \text{ lié } s'\\ \phantom{\text{b) }}(\text{i.e.\ } \{s,s'\}\in A).\end{array}\right.$

Pour l'instant on \uncertain{suppose} seulement l'axiome (1).

\emph{1 vérifié} Supposons $A\neq\emptyset$ (\uncertain{sinon} l'axiome est vide~!) donc (\add{en vertu de} l'axiome 1) $T\neq\emptyset$. Soit $t_0\in T$, $t=\{a,b,c\}$, avec $a,b,c\in S$.

Supposons de plus pour l'instant que (2) \emph{l'axiome 2 est vérifié pour $t_0$}, i.e.\ tout élément de $S\setminus t_0$ est lié à un des él.\ $a,b,c$ \add{(unique par l'axiome 1)}. Si $a\in t$, on désigne par $\widetilde{E}_a$ l'ens de $S\setminus t_0$ liés à $a$. Donc si $t=\{a,b,c\}$ on a bij.\ can.
\[
  S\setminus t_0=\coprod_{\alpha\in t_0}\widetilde{E}_{\alpha}=\widetilde{E}_a\amalg\widetilde{E}_b\amalg\widetilde{E}_c ,
\]
donc bij can
\[
  (1)\qquad S\simeq t_0\amalg\coprod_{\alpha\in t_0}\widetilde{E}_{\alpha}.
\]
\struck{D'autre part,}

\page{93}
Les arêtes $\subset t_0$ seront \uncertain{dites} horizontales \struck{\uncertain{basiques}}, celles qui lient des él.\ de $t_0$ à des él.\ de $S\setminus t_0$ seront \uncertain{dites} verticales, celles qui relient des pts de $S\setminus t_0$ \uncertain{seront dites} horizontales \uncertain{supérieures}. \struck{Il y a exactement} L'ens des arêtes \uncertain{basiques} \struck{\uncertain{Rappel}} est $\simeq\mathfrak{P}_2(t_0)$, l'ens des arêtes verticales est $\simeq\coprod_{\alpha\in t_0}\widetilde{E}_{\alpha}$ (\struck{\ill{}} à $s\in\widetilde{E}_{\alpha}$ correspond l'arête verticale $\{s,\alpha\}$).

Exprimons l'axiome 1 pour les arêtes envisagées jusqu'ici présent (celles qui ne sont pas hor.\ sup.) Pour une arête horizontale \uncertain{inférieure}, il est déjà exprimé par le fait que les $\widetilde{E}_{\alpha}$ ($\alpha\in t_0$) sont mutuellement disjoints, i.e.\ par l'écriture (1). Pour une arête verticale $\{\alpha,s\}$ ($\alpha\in t_0$, $s\in S\setminus t_0$), cela s'exprime ainsi :

\emph{Prop.~\uncertain{1}} \struck{\ill{}} $\forall\,\alpha\in t_0$, et $\forall\, s\in\widetilde{E}_{\alpha}$, $\exists!\, s'\in\widetilde{E}_{\alpha}$ lié à $s$.
\drawing{En marge, un triangle allongé de sommets $s$, $s'$ en haut et $a$ en bas ; plus bas, deux points $b$ et $c$.}

\emph{Cor.} $\alpha\to\alpha'$ est une involution sans pt fixe de $\widetilde{E}_{\alpha}$, notée $\sigma_{\alpha}$. On désigne par $\sigma$ l'application $S\to S$ telle que
\[
  \sigma|t_0=\mathrm{id}_{t_0},\qquad \sigma|\widetilde{E}_{\alpha}=\sigma_{\alpha}\quad\text{pour }\alpha\in t_0 .
\]
\note{Le numéro de cette proposition se lit plutôt « 2 » ; la page suivante renvoie à elle comme « Prop~1 » et numérote « Prop~2 » l'énoncé suivant. Le passage écrit aussi $\alpha\to\alpha'$ là où la proposition parle de $s\mapsto s'$ ; la lecture $\alpha\to\alpha'$ est celle de la page.}

\page{94}
Les arêtes \struck{\uncertain{supér}} horizontales supérieures de la forme $\{s,\sigma_{\alpha}(s)\}$ seront appelées supérieures internes (ou supérieures triviales) --- ce sont celles qui lient des pts d'un $=\widetilde{E}_{\alpha}$ ($\alpha\in t_0$). Exprimons l'axiome 1 pour ces nouvelles \struck{flèches} \add{arêtes}, on trouve

\emph{Prop.~2} Soient $s,s'\in\widetilde{E}_{\alpha}$ ($\alpha\in t_0$), $s$ et $s'$ liés (i.e.\ $s'=\sigma_{\alpha}s$). Soit $\beta\in t_0$, $\beta\neq\alpha$, et soit $t\in\widetilde{E}_{\beta}$. Alors $t$ ne peut être lié simultanément à $s$ et $s'$.

Appelons \struck{flèches} \add{arêtes} (\uncertain{hor.}) sup.\ transverses les \struck{flèches} \add{arêtes} (hor.) sup.\ qui ne sont pas internes, i.e.\ que \struck{ne sont} relient un pt d'un $\widetilde{E}_{\alpha}$ à un pt d'un $\widetilde{E}_{\beta}$ ($\alpha,\beta\in t_0$, $\alpha\neq\beta$). Soit $\{s,t\}$ ($s\in\widetilde{E}_{\alpha}$, $t\in\widetilde{E}_{\beta}$) une telle arête, \struck{\ill{}} considérons \add{$\{s,t,u\}$} le triangle \add{la contenant (\uncertain{il existe} par l'axiome 1)} \struck{correspondant (par l'axiome 1)}, mais \uncertain{pour} \uncertain{lequel} \uncertain{on se demande}~…). Évidemment $u\notin t_0$ (car \uncertain{sinon} …{} comme $u\in t_0$ on devrait avoir $u=\alpha$ et $u=\beta$, or $\alpha\neq\beta$~!) donc $u\in S\setminus t_0=\widetilde{E}_{\alpha}\amalg\widetilde{E}_{\beta}\amalg\widetilde{E}_{\gamma}$. On ne peut avoir $u\in\widetilde{E}_{\alpha}$, car on aurait alors $u=\sigma_{\alpha}(s)=s'$ par Prop~1, ce qui contredit prop.~2. Pour la même raison \struck{\ill{}} on a $u\notin\widetilde{E}_{\beta}$, donc $u\in\widetilde{E}_{\gamma}$, i.e.

\page{95}
\emph{Prop.~3} Si $t$ est un triangle contenant une arête sup.\ transverse (i.e.\ reliant un pt d'un $\widetilde{E}_{\alpha}$ et un d'un $\widetilde{E}_{\beta}$, \add{$\alpha,\beta\in t_0$, $\alpha\neq\beta$}) alors ses trois arêtes sont transverses, et la relation ``$\alpha\in t_0$ \uncertain{et} $\exists!\, s\in t$, $\alpha$ est lié à $s$ (\uncertain{avec} $s\in\widetilde{E}_{\alpha}$)'' est une bijection $t_0\xrightarrow{\sim}t$.

Les triangles en question s'appellent triangles supérieurs, leur ens est noté $\widetilde{T}$. On a donc des applications
\[
  \widetilde{T}\xrightarrow{\ \widetilde{\varphi}_{\alpha}\ }\widetilde{E}_{\alpha}\qquad(\alpha\in t)
\]
(on a

\begin{tikzcd}[column sep=small, row sep=small]
  \widetilde{E}_{a} & & \widetilde{E}_{b} \\
  & \widetilde{T} \arrow[ul, "\widetilde{\varphi}_{a}"] \arrow[ur, "\widetilde{\varphi}_{b}"'] \arrow[d, "\widetilde{\varphi}_{c}"] & \\
  & \widetilde{E}_{c} &
\end{tikzcd}

) \struck{définissant} \add{équivalentes à} une application unique
\[
  \widetilde{T}\xrightarrow{\ \widetilde{\varphi}\ }\prod_{\alpha\in t_0}\widetilde{E}_{\alpha}
  \quad\bigl(=\widetilde{E}_{a}\times\widetilde{E}_{b}\times\widetilde{E}_{c}\bigr).
\]
L'axiome 1 pour les arêtes sup.\ transverses équivaut donc à :

\emph{Prop.~4} Si \struck{$\alpha,\beta\in t_0$,} $\alpha=u\in\mathfrak{P}_2(t_0)$ est une

\page{96}
arête basique, \add{(disons $u=\{\alpha,\beta\}$)} l'application
\[
  \widetilde{T}\xrightarrow{\ \varphi_u\ }\prod_{i\in u}\widetilde{E}_i\quad\bigl(\simeq\widetilde{E}_{\alpha}\times\widetilde{E}_{\beta}\bigr)
\]
est \emph{injective}, et son image est formée des \struck{\uncertain{familles}} couples $(s,t)\in\widetilde{E}_{\alpha}\times\widetilde{E}_{\beta}$ tels que $s$ est lié à $t$.

\emph{Scholie} En résumé, moyennant les conditions (1) et (2) plus haut, on voit que la structure de graphe $S$, ainsi que les triangles $t_0$, se résument à la donnée de
\begin{itemize}
  \item[a)] un ens $t_0$ tel que $\mathrm{card}\,t_0=3$
  \item[b)] une famille $\widetilde{E}=(\widetilde{E}_{\alpha})_{\alpha\in t_0}$ d'ens.\ munis chacun d'une involution $\sigma_{\alpha}$ sans pt fixe
  \item[c)] une partie $\widetilde{T}\overset{\widetilde{\varphi}}{\hookrightarrow}\prod_{\alpha\in t_0}\widetilde{E}_{\alpha}$ \add{$=\widetilde{P}$}
\end{itemize}
[satisfaisant les conditions
\begin{itemize}
  \item[(*)] $\forall\, u\in\mathfrak{P}_2(t_0)$, l'application $\widetilde{T}\xrightarrow{\ \widetilde{\varphi}_u\ }\prod_{\alpha\in u}\widetilde{E}_{\alpha}$ est injection,
  \item[(**)] si $\alpha,\beta\in t_0$, $\alpha\neq\beta$, et si $s\in\widetilde{E}_{\alpha}$, $t\in\widetilde{E}_{\beta}$, si $(s,t)\in\mathrm{Im}\,\widetilde{\varphi}_{\{\alpha,\beta\}}(\widetilde{T})$, alors $(\sigma_{\alpha}s,t)\notin\mathrm{Im}\,\widetilde{\varphi}_{\{\alpha,\beta\}}(\widetilde{T})$.
\end{itemize}
\struck{et} de la façon suivante :]
\marginal{En travers de la marge gauche, encadré, une condition \uncertain{(***)} : « si $(s_{\alpha})_{\alpha\in t_0}\in\prod\widetilde{E}_{\alpha}$ et $(t_{\alpha})_{\alpha\in t_0}\in\widetilde{T}$ \uncertain{tels que} $s_{\beta}=\varphi_{\beta}(t_{\alpha})$ \ill{} $\forall\,\alpha,\beta$, $\alpha\neq\beta$, \ill{} » — rattachée par une accolade à (*) et (**).}
\note{La condition écrite dans la marge, tournée de 90°, n'est lue qu'en partie.}

1°) $S\simeq t_0\amalg\coprod_{\alpha\in t_0}\widetilde{E}_{\alpha}$

2°) \struck{deux pts} \struck{$\alpha$)} Deux pts de $t_0$ sont liés
\begin{itemize}
  \item[$\beta$)] un point de $\widetilde{E}_{\alpha}$ est lié à un et un seul él.\ de $t_0$, savoir $\alpha$ (si $\alpha\in t_0$)
  \item[$\gamma$)] un point $s$ de $\widetilde{E}_{\alpha}$ est lié à un \add{et un} seul pt de $\widetilde{E}_{\alpha}$, savoir $\sigma_{\alpha}(s)$
  \item[$\delta$)] \struck{Soit} \add{Soient} $\alpha,\beta\in t_0$, $\alpha\neq\beta$, $s\in\widetilde{E}_{\alpha}$, $t\in\widetilde{E}_{\beta}$, alors $s,t$ sont

\end{itemize}

\page{97}
liés ssi $(s,t)\in\mathrm{Im}\,\varphi_{\{\alpha,\beta\}}(\widetilde{T})$.

Étant donnés des données a) b) c), satisfaisant la condition (*), le graphe qu'on vient de décrire satisfait les conditions (1) (Axiome 1) et (2) (L'axiome 2 \add{seulement} pour le triangle $t_0$).

L'ens $T$ des triangles de $(S,A)$ est alors donné, à une bijection près, par
\[
  T\simeq\underbrace{\{t_0\}}_{\substack{\text{triangle}\\ \text{de base}}}\amalg\underbrace{\coprod_{\alpha\in t_0}E_{\alpha}}_{\substack{\text{triangles}\\ \text{verticaux}}}\amalg\underbrace{\widetilde{T}}_{\substack{\text{triangles (horiz.)}\\ \text{supérieurs}}}
\]
où on pose \struck{\ill{}}
\[
  E_{\alpha}=\widetilde{E}_{\alpha}/\sigma_{\alpha}
\]
(de sorte que $\widetilde{E}_{\alpha}$ est un revêt.\ de degré 2 de $E_{\alpha}$). Pour $\alpha\in t_0$, un élément $u$ de $E_{\alpha}$, i.e.\ $u\in\mathfrak{P}_2(\widetilde{E}_{\alpha})$ formé d'une orbite sous $\sigma_{\alpha}$, \add{(vertical)} le triangle associé est par définition $\{u\cup\{\alpha\}\}$. Pour \struck{\ill{}} $\tau\in\widetilde{T}$, le triangle associé est par définition l'ens $\{\varphi_{\alpha}(\tau)\}_{\alpha\in t_0}$.

Rien n'empêche pour l'instant qu'on ait $\widetilde{T}=\emptyset$~! Mais \uncertain{(on va)} \uncertain{proposer} maintenant d'exprimer l'axiome 2 pour \emph{tous} les triangles de $S$ (pas seulement $t_0$). Prenons d'abord les

\page{98}
Cas des triangles \emph{verticaux} $\{\alpha,s,s'\}$ ($\alpha\in t_0$, $s\in\widetilde{E}_{\alpha}$, $s'=\sigma_{\alpha}(s)$). On trouve

\emph{Prop.~5} Soient $\alpha\in t_0$, $s,s'$ deux arêtes liées de $\widetilde{E}_{\alpha}$, et \struck{$\beta\in t_0$, $\beta\neq\alpha$,} $t\in\widetilde{E}_{\beta}$, où $\beta\in t_0$, $\beta\neq\alpha$, alors $t$ est lié à \struck{au plus} $s$ ou à $s'$ (\emph{NB} par prop.~2, on \uncertain{ne peut} \uncertain{avoir} les deux \struck{il est lié à un} \uncertain{à la fois}, il \uncertain{viendrait} lié à $s$ et à $s'$, donc il est lié à un et un seul des \emph{deux} pts $s,s'$). En termes des données de la Scholie : Soit $\widetilde{T}_{\alpha,\beta}\subset\widetilde{E}_{\alpha}\times\widetilde{E}_{\beta}$ l'image de $\widetilde{T}$ dans $\widetilde{E}_{\alpha}\times\widetilde{E}_{\beta}$, soient $s,s'\in\widetilde{E}_{\alpha}$, $t\in\widetilde{E}_{\beta}$, avec $s'=\sigma_{\alpha}s$ i.e.\ $s=\sigma_{\alpha}s'$. Alors on a $\{s,t\}\in\widetilde{T}_{\alpha\beta}$ ou $\{s',t\}\in\widetilde{T}_{\alpha\beta}$.

\emph{Cor.} Sous les conditions précédentes, posant $t'=\sigma_{\beta}t$, \add{$s,t$ liés i.e.} si \struck{$(s,t)\in\widetilde{T}$} $\{s,t\}\in A$ i.e.\ $(s,t)\in\widetilde{T}_{\alpha\beta}$, alors \add{$s',t'$ liés i.e.} $\{s',t'\}\in A$ i.e.\ $(s',t')\in\widetilde{T}_{\alpha\beta}$.

En effet, $t'$ étant lié à $s$ ou $s'$, et ne pouvant être lié à $s$ (car $s$ étant lié à $t$ et $t'$ \uncertain{contredirait} prop.~2) il est lié à $s'$.

\emph{Cor.} Soit $\widetilde{\sigma}_{\widetilde{P}}=\prod_{\alpha\in t_0}\sigma_{\alpha}$ l'involution \uncertain{(sans pt fixe)} d'ordre 2 \struck{\ill{}} de $\widetilde{P}=\prod\widetilde{E}_{\alpha}$. Alors $\widetilde{T}$ est stable par $\sigma_{\widetilde{P}}$\,; on note $\sigma_{\widetilde{T}}$ l'automorphisme d'ordre 2 de $\widetilde{T}$ induit par $\sigma_{\widetilde{P}}$.

\page{99}
\emph{Cor.} Soient $\alpha\in t_0$, $s\in\widetilde{E}_{\alpha}$, \struck{$\beta,\gamma\in t_0\setminus\{\alpha\}$, $\beta\neq\gamma$} \add{$\beta\in t_0\setminus\{\alpha\}$}, soit $\widetilde{E}_{\beta}(s)$ \add{$\subset\widetilde{E}_{\beta}$} l'ens des $t\in\widetilde{E}_{\beta}$ liés à $s$, c'est \add{donc} une \emph{section} de $\widetilde{E}_{\beta}$ sur $E_{\beta}=\widetilde{E}_{\beta}/\sigma_{\beta}$. Soit \add{$\gamma\in t_0\setminus\{\alpha,\beta\}$} \struck{$\gamma\in t_0$, $\gamma\neq\beta$}, d'où $\widetilde{E}_{\gamma}(s)$ section de $\widetilde{E}_{\gamma}$ sur $E_{\gamma}$\,; la relation entre $t\in\widetilde{E}_{\beta}(s)$ et $t'\in\widetilde{E}_{\gamma}(s)$ \,``$t$ et $t'$ sont liés i.e.\ $(t,t')\in\widetilde{T}_{\beta,\gamma}$'' est une bijection \struck{\ill{}}
\[
  u_s^{\gamma,\beta}:\widetilde{E}_{\beta}(s)\xrightarrow{\ \sim\ }\widetilde{E}_{\gamma}(s).
\]

\emph{Cor.} \struck{Si} \uncertain{Si} $\widetilde{E}_{\alpha}\neq\emptyset$, alors $\mathrm{card}\,\widetilde{E}_{\beta}=\mathrm{card}\,\widetilde{E}_{\gamma}$ donc $\mathrm{card}\,\widetilde{E}_{\beta}=\mathrm{card}\,\widetilde{E}_{\gamma}$\,; si on \struck{\ill{}} cardinal est $\neq0$ i.e.\ $\widetilde{E}_{\beta}\neq\emptyset$ ($\Longleftrightarrow\widetilde{E}_{\gamma}\neq\emptyset$) alors $\widetilde{E}_{\alpha},\widetilde{E}_{\beta},\widetilde{E}_{\gamma}$ ont \uncertain{même} cardinal (et $E_{\alpha},E_{\beta},E_{\gamma}$ ont \uncertain{même} cardinal).

\emph{NB} Évidemment, si $\widetilde{E}_{\alpha}=\emptyset$ $\forall\,\alpha\in t_0$, les axiomes \struck{\ill{}} 1 et 2 sont \uncertain{dans ce cas} satisfaits. De même si $\widetilde{E}_{\alpha}\neq\emptyset$, et $\widetilde{E}_{\beta},\widetilde{E}_{\gamma}=\emptyset$ (donc $\widetilde{T}=\emptyset$, donc pas d'\struck{flèches} arêtes sup transverses). Alors le \struck{\ill{}} \add{graphe} est \struck{\ill{}} de la forme
\[
  S=\{\alpha\}\amalg\widetilde{S}'\qquad\text{où }\widetilde{S}'=S-\{\alpha\}
\]
est muni d'une invol.\ $\sigma_{S'}$ sans pt fixe, et les arêtes sont les $\{\alpha,s\}$ ($s\in\widetilde{S}'$) et les $(s,\sigma_{S'}s)$ ($s\in\widetilde{S}'$), donc
\[
  A\simeq\widetilde{S}'\amalg\underbrace{\widetilde{S}'/\sigma_{S'}}_{S'} .
\]
\struck{Mais}

\page{100}
On supposera désormais que \struck{l'axiome}

\emph{Axiome 3} Les $\widetilde{E}_{\alpha}$ ($\alpha\in t_0$) sont $\neq\emptyset$, i.e.\ \uncertain{on est} \add{(trivial)} pas dans le cas (envisagé dans le NB précédent) \uncertain{où} \uncertain{un} sommet $s\in S$ \uncertain{est} \uncertain{dans} au moins 2 triangles.

Je veux exprimer, dans ce cas, l'axiome 2 pour les triangles sup. Or il est \uncertain{automatique}~!
\note{Le passage qui suit (Prop.~6 et sa démonstration, jusqu'à « c'est absurde ») est barré de trois longs traits obliques au crayon ; il reste lisible et est transcrit en entier.}

\drawing{Dans la marge : un prisme triangulaire en perspective, triangle supérieur $\{t,s,u\}$, triangle inférieur $\{b,a,c\}$ (avec $a$ au fond), arêtes verticales.}

\emph{Prop.~6} Soit $\{s,t,u\}$ \uncertain{un} $\{a,b,c\}=t_0$ un triangle supérieur (i.e.\ un $\in\widetilde{T}$), et soit $s_1\in\widetilde{E}_a\setminus\{s,s'=\sigma_a s\}$. Alors $s_1$ est lié à $t$ ou $u$ (\emph{NB} bien sûr, pas aux deux à la fois…)

\emph{Démonstration} : \struck{En effet,} \struck{\ill{}} a) \struck{tout} $\alpha\in t_0$ est lié à un pt de $\{s,t,u\}$ b) Si \struck{$s_1\in\widetilde{E}_{\alpha}$, alors} $S\setminus t_0=\coprod\widetilde{E}_{\alpha}$, \struck{s distinct de} disons $s_1\in E_a$, \struck{\ill{} $\alpha\in t_0$ \ill{} $s_1$} $\notin t$ i.e.\ $s_1\neq s$, si $s_1$ \struck{$=\sigma_a s=s'$, il est} \ill{} n'est pas lié à $s$ i.e.\ $s_1\neq s'=\sigma_a s$, montrons qu'il est lié à $t$ ou $u$. En effet, s'il n'était lié à aucun des deux, alors $s_1'=\sigma_a s_1$ serait lié aux deux : il \uncertain{suit}, et comme $s_1'\neq s$ c'est absurde.

\emph{Résolution} En résumé, la donnée d'un complexe $(S,A)$ \struck{et \ill{}} satisfaisant les axiomes (1), (2) et \uncertain{que} …{} (\uncertain{discret}, \uncertain{ou réduit à un triangle, ou à un} …{} du type trivial \struck{étoile} \uncertain{étoile} de triangles
\drawing{Une étoile de triangles : plusieurs triangles ayant un sommet commun, disposés en roue.}
) équivaut à celle de
\note{La phrase se poursuit au-delà de ce lot.}

\end{document}
